\chapter{Results}

\section{Introduction}

This chapter reports the compatibility, architecture, ablation, overfitting,
bias and candidate-search results. Results are separated according to their
experimental protocol. In particular, the selected compatibility model, the
three-seed architecture comparison and the candidate-search experiment are
reported independently because they answer different questions.

\section{Selected compatibility-model performance}

The selected compatibility model was evaluated on 16,062 held-out test
examples containing equal numbers of positive and negative rows. It achieved a
test ROC--AUC of 0.8583, with a 95\% confidence interval from 0.8524 to 0.8641.
Accuracy and balanced accuracy were both 0.7664. \Cref{tab:final-model-results}
reports the remaining classification results.

\begin{table}[htbp]
  \centering
  \caption{Held-out performance of the selected compatibility model}
  \label{tab:final-model-results}
  \begin{tabular}{lr}
    \toprule
    Metric & Result \\
    \midrule
    ROC--AUC & 0.8583 \\
    95\% AUC confidence interval & 0.8524--0.8641 \\
    Average precision & 0.8558 \\
    Accuracy & 0.7664 \\
    Balanced accuracy & 0.7664 \\
    Precision & 0.7240 \\
    Recall & 0.8610 \\
    F1 & 0.7866 \\
    Specificity & 0.6718 \\
    Matthews correlation coefficient & 0.5426 \\
    Brier score & 0.1599 \\
    \bottomrule
  \end{tabular}
\end{table}

At the 0.5 threshold, the model correctly classified 5,395 negative and 6,915
positive outfits. It classified 2,636 negative outfits as positive and 1,116
positive outfits as negative. Recall was therefore higher than specificity at
this threshold.

A separate fixed-seed baseline experiment produced an AUC of 0.4940 for
deterministic random scores and 0.7012 for mean pairwise SigLIP cosine
similarity. The learned head in that experiment reached 0.8573. These values
came from an earlier controlled baseline run and are reported separately from
the selected model result of 0.8583. A mean-pooled frozen-SigLIP logistic
baseline, evaluated on the same processed test rows, reached AUC 0.5532. Its
deterministic preprocessing produced the same value in the three recorded seed
files.

\section{Controlled architecture comparison}

Each learned architecture was trained with seeds 42, 7 and 123. The proposed
model achieved the highest mean test AUC and the smallest AUC standard deviation
across the three runs (\Cref{tab:architecture-results}).

\begin{table}[htbp]
  \centering
  \caption{Controlled three-seed architecture comparison}
  \label{tab:architecture-results}
  \begin{tabularx}{\textwidth}{>{\raggedright\arraybackslash}Xrrrr}
    \toprule
    Model & Parameters & Time (s) & Mean AUC & SD \\
    \midrule
    Bidirectional LSTM & 397,185 & 96.87 & 0.7928 & 0.0060 \\
    Type-aware pairwise & 591,489 & 7.99 & 0.7918 & 0.0129 \\
    Set Transformer & 374,145 & 7.03 & 0.7744 & 0.0133 \\
    Proposed model & 379,265 & 4.81 & 0.8596 & 0.0009 \\
    \bottomrule
  \end{tabularx}
\end{table}

The proposed model's mean balanced accuracy was 0.7742 and its mean Matthews
correlation coefficient was 0.5510. Mean balanced accuracy was 0.7138 for the
bidirectional LSTM, 0.7176 for the type-aware model and 0.7034 for the set
Transformer.

Paired bootstrap comparison of the three-seed mean predictions gave an AUC
difference of 0.0700 between the proposed model and the bidirectional LSTM
(95\% interval 0.0650--0.0753). The difference from the type-aware model was
0.0634 (0.0578--0.0694), while the difference from the set Transformer was
0.0801 (0.0750--0.0858). These results compare controlled implementations on
the project's processed data rather than the headline results in the original
publications.

\section{Fill-in-the-blank and ablation results}

Of 15,145 official disjoint FITB questions, 4,468 were compatible with the
project's four-position representation, giving coverage of 29.50\%. Other
questions were excluded because of unsupported answer positions, insufficient
supported garments, duplicate positions or an answer position already occupied
in the question.

The selected model achieved FITB accuracy of 0.6025. Its 95\% bootstrap
confidence interval was 0.5873--0.6175, compared with chance accuracy of 0.25.
Mean reciprocal rank was 0.7705 and normalised discounted cumulative gain was
0.8289.

The structural ablation used the same 4,468 questions and three random seeds.
The 0.6025 result above belongs to the selected seed-42 checkpoint; the complete
model row below is the mean of three separately retrained seed runs and is
therefore 0.5994.
\Cref{tab:ablation-results} reports the mean accuracy and the paired difference
between the complete model and each reduced variant.

\begin{table}[htbp]
  \centering
  \caption{Structural ablation on the retained official FITB subset}
  \label{tab:ablation-results}
  \begin{tabularx}{\textwidth}{>{\raggedright\arraybackslash}Xrrr}
    \toprule
    Configuration & Mean accuracy & Full minus variant & Paired 95\% interval \\
    \midrule
    Complete model & 0.5994 & --- & --- \\
    Without clothing-position embedding & 0.5947 & 0.0047 & -0.0013--0.0106 \\
    Without pairwise module & 0.4966 & 0.1028 & 0.0896--0.1157 \\
    \bottomrule
  \end{tabularx}
\end{table}

The retained FITB data contained only 46 menswear questions, compared with
4,422 womenswear questions. Menswear performance was therefore recorded as a
coverage limitation rather than treated as a reliable comparison between
clothing ranges. The descriptive subgroup values are provided in the
appendices.

\section{Calibration and overfitting audit}

The selected checkpoint reached its highest validation AUC at epoch 6. Training,
validation and test AUC were 0.9699, 0.8598 and 0.8583 respectively
(\Cref{tab:split-results}).

\begin{table}[htbp]
  \centering
  \caption{Performance of the selected checkpoint by data split}
  \label{tab:split-results}
  \begin{tabular}{lrrrr}
    \toprule
    Split & Examples & AUC & Accuracy & F1 \\
    \midrule
    Training & 18,243 & 0.9699 & 0.8997 & 0.9051 \\
    Validation & 3,274 & 0.8598 & 0.7688 & 0.7845 \\
    Test & 16,062 & 0.8583 & 0.7664 & 0.7866 \\
    \bottomrule
  \end{tabular}
\end{table}

The train--test AUC difference was 0.1116, while the absolute
validation--test difference was 0.00145. Validation AUC declined by 0.0100
between the saved epoch and the last observed epoch.

Validation-only temperature scaling produced a temperature of 1.5760.
Validation negative log-likelihood decreased from 0.5081 to 0.4776, while
expected calibration error decreased from 0.0772 to 0.0588. Candidate ordering
was unchanged by the calibration step. Expected calibration error used ten
equal-width probability bins over the 3,274 validation examples.

\section{Outfit-length imbalance and feature-domain results}

The class labels were balanced. A separate experiment examined the imbalance
between three-item and four-item outfits. \Cref{tab:length-weighting-results}
reports the three configurations across the same three seeds.

\begin{table}[htbp]
  \centering
  \caption{Outfit-length weighting experiment}
  \label{tab:length-weighting-results}
  \begin{tabularx}{\textwidth}{>{\raggedright\arraybackslash}Xrrr}
    \toprule
    Configuration & Overall AUC & Weakest-length AUC & Group gap \\
    \midrule
    Unweighted & 0.8596 & 0.8507 & 0.0393 \\
    Moderate length weighting & 0.8601 & 0.8500 & 0.0433 \\
    Stronger length weighting & 0.8553 & 0.8438 & 0.0485 \\
    \bottomrule
  \end{tabularx}
\end{table}

Moderate weighting achieved the highest overall mean AUC but reduced the mean
result for the weakest length group and increased the group gap. It therefore
failed the predefined deployment gate, and the unweighted configuration was
retained.

The feature-domain diagnostic compared 26,124 Polyvore test items with 300 Zara
catalogue items. The domain classifier achieved cross-validated AUC of 1.000.
Centroid cosine similarity was 0.7452, and the mean nearest-Polyvore similarity
for the Zara items was 0.6528. The implications of this separability are
considered in Chapter 5.

\section{Candidate-search results}

The product-independent candidate artefact contained 1,670 clothing concepts
derived from the Polyvore training split, covering 23 garment types, 12 colours
and nine styles. The selected womenswear range contained 1,508 concepts and the
menswear range contained 162. Both ranges contained candidates for all four
supported clothing positions.

The earlier retailer-derived candidate space contained 300 product records and
122 visible type--colour combinations. The product-independent representation
contained 271 such combinations. Absolute compatibility scores across these
candidate sources were not treated as directly comparable because their data
distributions differ.

Exact and quota-limited search were compared on 18 held-out inputs using the
same product-independent candidate space and compatibility model. The
quota-limited path recovered the exact top result in 27.78\% of queries and had
mean recall of 24.44\% for the exact top ten. It examined a mean 3.41\% of the
candidate combinations.

The mean score loss relative to the highest score found by exact search was
0.0179, with a 95\% bootstrap interval of 0.0063--0.0357. Thirteen of the 18
queries had a positive score loss. Repeated exact searches returned the same
top-ten results for every tested query.

\begin{table}[htbp]
  \centering
  \caption{Exact-search and quota-limited search results}
  \label{tab:search-results}
  \begin{tabular}{lr}
    \toprule
    Measure & Quota-limited result \\
    \midrule
    Exact top result recovered & 27.78\% \\
    Mean recall of exact top ten & 24.44\% \\
    Mean candidate space examined & 3.41\% \\
    Mean score loss & 0.0179 \\
    Queries with positive score loss & 13 of 18 \\
    \bottomrule
  \end{tabular}
\end{table}

On the recorded CPU experiment, exact search had a median runtime of 0.474
seconds, a 95th-percentile runtime of 2.107 seconds and a maximum of 3.206
seconds. Quota-limited search over the same candidate space had a median runtime
of 0.028 seconds.

Candidate availability was tested across four weather scenarios and nine
styles, giving 36 combinations. Every tested combination retained all four
logical clothing positions. In hot-weather cases, the outer-layer position was
represented by the explicit masked \emph{No Outer Layer} option.

\section{Summary}

The selected compatibility model achieved test AUC of 0.8583 and FITB accuracy
of 0.6025 on the retained official subset. The three-seed comparison recorded a
higher mean AUC for the proposed architecture than for the three controlled
learned baselines. Removing the pairwise module produced a clear reduction,
whereas the interval for removing clothing-position information crossed zero.

The selected model showed a substantial train--test gap but close validation
and test performance. Validation-only calibration reduced negative
log-likelihood and expected calibration error. The tested outfit-length
weighting did not pass its deployment gate.

The product-independent candidate representation expanded the range of visible
type--colour combinations. The quota-limited search recovered the same highest
scoring result as exact search in 27.78\% of the tested queries while requiring
substantially less runtime. Chapter 5 examines the meaning and limitations of
these results.
